%!TEX root = ../../main.tex
%% ===============================================================
%% 1.5 논문의 구성
%% 파일: chapters/01_introduction/05_outline.tex
%% 상위: chapters/01_introduction/chapter.tex
%% 정의 라벨: sec:intro-outline
%% 참조 라벨: ch:background, ch:conclusion, ch:discussion, ch:features, ch:label, ch:localization, ch:models, ch:problem, ch:results, ch:setup
%% 포함 객체: -
%% 인용: -
%% 상태: 초안 (docs/OUTLINE.md 참고)
%% ===============================================================

\section{논문의 구성}
\label{sec:intro-outline}

제 \ref{ch:background} 장은 관련 연구와 본 논문이 사용하는 개념들의 정의·기원을 정리한다. 제 \ref{ch:problem} 장은 텔레메트리와 교전, 예측 문제를 형식적으로 정의한다. 제 \ref{ch:localization} 장은 교전 국소화 알고리즘과 그 타당성 검증 설계를, 제 \ref{ch:label} 장은 교환가치 라벨을, 제 \ref{ch:features} 장은 특징 표현을 다룬다. 제 \ref{ch:models} 장은 비교 대상 모델과 일곱 가지 처치를 정의하고, 제 \ref{ch:setup} 장은 데이터·분할·학습·평가 절차를, 제 \ref{ch:results} 장은 결과를 제시한다. 제 \ref{ch:discussion} 장에서 설계 결정, 코드--논문 감사, 타당성 위협을 논의하고, 제 \ref{ch:conclusion} 장에서 결론과 향후 연구를 제시한다.
